\section{实用工具}

\subsection{Windows 自动编译器}

\lstinputlisting[language=bash]{codes/07-utilities-01.bat}

\textbf{使用方法}
\lstinputlisting[language=bash]{codes/07-utilities-02.bat}

\subsection{随机数生成器}

\lstinputlisting[language=C++]{codes/07-utilities-03.cpp}

\subsection{Python 随机数生成}

\textbf{用途：} 使用 \code{random} 标准库为对拍生成随机数据，或为随机化算法提供随机源。

\begin{itemize}
  \item \code{random.seed(x)}：固定随机种子，使生成结果可复现；不传参则由系统熵源自动播种。
  \item \code{random} 模块是\textbf{确定性伪随机}（梅森旋转算法），输出序列可被预测；对抗 hack 的场景应改用 \code{random.SystemRandom()}。
  \item 批量生成大量随机数时，\code{random.choices} / \code{random.sample} 由 C 实现，比循环调用 \code{randint} 快数倍。
\end{itemize}

\lstinputlisting[language=Python]{codes/07-utilities-09.py}

\subsection{随机树图生成器}

全随机树图：
\lstinputlisting[language=C++]{codes/07-utilities-04.cpp}

生成单链：
\lstinputlisting[language=C++]{codes/07-utilities-05.cpp}

\subsection{计时器}

\lstinputlisting[language=C++]{codes/07-utilities-06.cpp}

\subsection{单文件对拍}

\lstinputlisting[language=C++]{codes/07-utilities-07.cpp}

\subsection{Python 对拍器}

用于对拍两个程序的输出，需要准备三个可执行文件：
\begin{itemize}
    \item \code{gen}：随机数据生成器
    \item \code{my}：你的待测程序
    \item \code{std}：标准/暴力程序
\end{itemize}

\lstinputlisting[language=Python]{codes/07-utilities-08.py}
